% ID: SEL-D0064-OB-E02
% ESTADO: validada. S064.
\textbf{a)} Para continuidad en $1$ igualamos $f(1)=2-2a$ al límite derecho $a-1$. Resulta $\boxed{a=1}$.

\textbf{b)} Con $a=1$, los tramos son $4-(x+1)^2$ para $x\le1$ y $(x-3)^2-4$ para $x>1$. La primera parábola tiene vértice $(-1,4)$ y la segunda $(3,-4)$. Se unen en $(1,0)$.
\begin{center}
\setlength{\unitlength}{0.9mm}
\begin{picture}(112,67)(-14,-9)
\qbezier(-7,30)(39,30)(85,30)
\qbezier(30,1)(30,27)(30,53)
\qbezier(0,30)(20,70)(40,30)
\qbezier(40,30)(60,-10)(80,30)
\put(86,30){$x$}\put(32,53){$f(x)$}
\put(13,52){\scriptsize $(-1,4)$}
\put(38,32){\scriptsize $(1,0)$}
\put(55,4){\scriptsize $(3,-4)$}
\end{picture}
\end{center}
La gráfica continúa hacia la izquierda por la primera parábola y hacia la derecha por la segunda. Crece en $(-\infty,-1]$ y $[3,\infty)$ y decrece en $[-1,3]$. El máximo local es $(-1,4)$ y el mínimo local es $(3,-4)$. En $x=1$ el descenso continúa, así que no hay extremo.
